% ---
% Inicia os anexos
% ---
\begin{anexosenv}
    \chapter{Código Questao 3.4}\label{cap:codigo_q34}
    
    \begin{lstlisting}[language=Python, breaklines=true, basicstyle=\tiny]
    import numpy as np
    import matplotlib.pyplot as plt
    import pandas as pd

    def f(x):
        if 8 <= x <= 10: return x - 8
        elif 10 < x <= 12: return 12 - x
        else: return 0  

    # Parameters
    dt = 5e-12 
    dx = 0.1
    x_max = 20 # transmission line length
    t_max = 2e-9 # final time
    L = 4e-9
    C = 1.6e-12
    Z0 = np.sqrt(L / C) # characteristic impedance
    vp = 1 / np.sqrt(L * C) # wave velocity

    # Parametros de grade
    nx = int(x_max / dx) + 1
    nt = int(t_max / dt) + 1
    C2 = (vp * dt / dx)**2 

    # Inicializacao da matriz de potencial V[tempo, espaco]
    V = np.zeros((nt, nx))

    # condicao inicial
    x = np.linspace(0, x_max, nx)
    initial = np.zeros(nx)

    for i in range (nx):
        initial[i] = f(x[i])

    V[0, :] = initial

    # Primeiro passo de tempo
    for i in range(1, nx - 1):
        V[1, i] = V[0, i]

    # Loop Principal
    for n in range(1, nt - 1):
        V[n+1, 1:-1] = C2 * (V[n, 2:] - 2*V[n, 1:-1] + V[n, :-2]) + 2*V[n, 1:-1] - V[n-1, 1:-1]

    fig, axes = plt.subplots(nrows=3, ncols=3, figsize=(10, 10))

    t = np.linspace(0, t_max, nt)

    t_list = [0, 10, 30, 40, 50, 60, 70, 80, 90]

    for i, ax in enumerate(axes.flatten()):
        ax.plot(x, V[t_list[i], :], color='blue') # Plot the i-th dataset on the i-th axis
        ax.set_title(f'Instante {t_list[i] * dt * 1e9:.2f} ns')      # Set a title for each subplot
        ax.set_xlabel('Posicao (cm)')
        ax.set_ylabel('Tensao (V)')
        ax.grid(True)

    plt.tight_layout(rect=[0, 0.03, 1, 0.95]) # Adjust rect to make space for suptitle

    plt.show()
    \end{lstlisting}
    
\end{anexosenv}